Antes de demostrar, recordemos qué hace cada modelo. Ambos reciben un vector
$\mathbf{x} \in \mathbb{R}^d$ y producen una etiqueta $\hat{y} \in \{0, 1\}$.
El perceptrón lo hace comparando la combinación lineal $z = \mathbf{w}^\top \mathbf{x} + b$
directamente contra cero mediante una función escalón. La regresión logística, en cambio,
primero transforma $z$ en una probabilidad a través de la función sigmoide
\[
\sigma(z) = \frac{1}{1+e^{-z}},
\]
que es estrictamente creciente y toma valores en $(0,1)$, y luego compara esa
probabilidad contra el umbral $0.5$. La clave de la demostración es que $\sigma$,
al ser estrictamente creciente, preserva el orden: $\sigma(z) \ge 0.5$ ocurre
exactamente cuando $z \ge 0$, de modo que ambas comparaciones son equivalentes.\\

\textbf{Demostración.} Sea un vector de entrada $\mathbf{x}$ y supongamos que tanto
el perceptrón como la regresión logística utilizan los mismos parámetros $\mathbf{w}$
y $b$. En ambos modelos se calcula primero la combinación lineal
\[
z = \mathbf{w}^T \mathbf{x} + b.
\]

\textbf{Perceptrón.} El perceptrón clasifica mediante la función escalón:
\[
\hat{y} =
\begin{cases}
1, & z \ge 0, \\
0, & z < 0.
\end{cases}
\]
Por tanto, su regla de decisión es
\[
\hat{y} =
\begin{cases}
1, & \mathbf{w}^T\mathbf{x} + b \ge 0, \\
0, & \mathbf{w}^T\mathbf{x} + b < 0.
\end{cases}
\]

\textbf{Regresión logística.} La regresión logística calcula la probabilidad
\[
P(y=1 \mid \mathbf{x}) = \sigma(z) = \frac{1}{1+e^{-z}}.
\]
La clasificación se obtiene utilizando el umbral $0.5$:
\[
\hat{y} =
\begin{cases}
1, & \sigma(z) \ge 0.5, \\
0, & \sigma(z) < 0.5.
\end{cases}
\]

Ahora demostramos que $\sigma(z) \ge 0.5 \iff z \ge 0$. En efecto,
\[
\frac{1}{1+e^{-z}} \ge \frac{1}{2}
\iff
2 \ge 1 + e^{-z}
\iff
1 \ge e^{-z}
\iff
0 \ge -z
\iff
z \ge 0.
\]
Por consiguiente, $\sigma(z) \ge 0.5 \iff z \ge 0$. Sustituyendo
$z = \mathbf{w}^T\mathbf{x} + b$, la regla de clasificación de la regresión
logística queda
\[
\hat{y} =
\begin{cases}
1, & \mathbf{w}^T\mathbf{x} + b \ge 0, \\
0, & \mathbf{w}^T\mathbf{x} + b < 0.
\end{cases}
\]
Esta regla coincide exactamente con la del perceptrón. Por tanto,
\[
\hat{y}_{\text{Perceptrón}} = \hat{y}_{\text{Logística}}.
\]
En conclusión, bajo un mismo conjunto de parámetros $\mathbf{w}$ y $b$, el perceptrón
y la regresión logística inducen la misma frontera de decisión $\{\mathbf{x} :
\mathbf{w}^T\mathbf{x} + b = 0\}$ y realizan exactamente la misma clasificación
para todo punto $\mathbf{x} \in \mathbb{R}^d$. \hfill$\blacksquare$
